\chapter{Appendices}

\section{Local Development Commands}

The repository includes scripts for both Docker-based and local development. The Docker path starts the complete service stack:

\begin{verbatim}
cd /Users/omer/Desktop/bitirme
bash scripts/docker_up.sh
\end{verbatim}

For local backend development:

\begin{verbatim}
cd backend
python3 -m venv .venv
source .venv/bin/activate
pip install -r requirements.txt
uvicorn app.main:app --reload
\end{verbatim}

For local frontend development:

\begin{verbatim}
cd frontend
npm install
npm run dev
\end{verbatim}

\section{Important API Endpoints}

\begin{table}[!htbp]
    \centering
    \caption{\label{tab:api-endpoints}Representative API endpoints.}
    \begin{tabularx}{\textwidth}{lX}
        \toprule
        Endpoint & Purpose \\
        \midrule
        \texttt{GET /api/health} & Health check. \\
        \texttt{POST /api/questions/manual} & Ask and answer a manual question immediately. \\
        \texttt{POST /api/questions/manual/queue} & Add a manual question to the live queue. \\
        \texttt{GET /api/live/state} & Return current broadcast state for the live stage. \\
        \texttt{POST /api/admin/archive/crawl} & Download and archive official sources. \\
        \texttt{POST /api/admin/archive/index} & Index archived sources into documents and chunks. \\
        \texttt{GET /api/admin/archive/stats} & Return archive statistics. \\
        \texttt{POST /api/streams/youtube/connect} & Connect a YouTube live chat session. \\
        \bottomrule
    \end{tabularx}
\end{table}

\section{Evaluation Cases}

The quality suite uses 12 questions. Table \ref{tab:evaluation-cases} gives a compact overview.

\begin{table}[!htbp]
    \centering
    \caption{\label{tab:evaluation-cases}Evaluation question categories.}
    \begin{tabularx}{\textwidth}{llX}
        \toprule
        Category & Count & Examples \\
        \midrule
        Regulation & 4 & undergraduate duration, graduate regulation, thesis guide, diploma directive \\
        Web & 3 & ECTS package, academic calendar, regulation page \\
        Department & 2 & Computer Engineering mission and target student profile \\
        PDF & 3 & research policy, Erasmus directive, student guide \\
        \bottomrule
    \end{tabularx}
\end{table}

\section{Environment Variables}

The main environment variables are:

\begin{itemize}
    \item \texttt{AI\_PROVIDER}: selects \texttt{openai} or \texttt{openrouter};
    \item \texttt{OPENAI\_API\_KEY} and \texttt{OPENROUTER\_API\_KEY}: provider keys;
    \item \texttt{AI\_CHAT\_MODEL} and \texttt{AI\_EMBEDDING\_MODEL}: optional model overrides;
    \item \texttt{YOUTUBE\_API\_KEY}: key used for live chat polling;
    \item \texttt{DATABASE\_URL}: database connection string;
    \item \texttt{TTS\_ENABLED}: runtime default for speech generation.
\end{itemize}
